\acterecto{Intention d’indulgences}
\usage{quot., au lever}
\ouverture{M}{on}{Dieu,} je Vous offre toutes les œuvres de cette journée, dans l’intention de Vous glorifier, de réparer mes péchés et de satisfaire pour les âmes du purgatoire. Je désire gagner, par Votre grâce, toutes les indulgences attachées à mes actions de ce jour, autant qu’il Vous plaira de m’en accorder. Faites que rien, en moi, ne Vous en empêche.
\culdelampe

\acteverso{Acte de soumission à la volonté de Dieu}
\usage{quot., le matin}
\ouverture{M}{on}{Dieu,} je reconnais que Vos vues sont très supérieures aux miennes. Aussi tout ce que Vous permettrez aujourd’hui, dans l’ordre du corps, de l’âme, des circonstances ou des épreuves, je l’accepte d’avance, pleinement, volontairement, sans réserve, comme ce qu’il y a de meilleur pour mon salut et pour Votre gloire.

Je consens aussi, ô mon Dieu, à être ignoré, incompris, contredit, oublié, méprisé et calomnié si cela Vous glorifie et m’abaisse. Ne permettez pas que je me défende autrement qu’en me taisant.

Je Vous remets, Seigneur, ma volonté propre, mes préférences, mes attentes, mes inquiétudes, et je Vous prie humblement de faire en moi ce qu’il Vous plaira.

Et parce que je suis faible, Seigneur, et prompt à me plaindre quand la douleur ou la contradiction surviennent, je Vous demande de ne pas tenir compte de mes murmures, de mes impatiences ou de mes révoltes naturelles. Voyez-y non l’opposition de ma volonté, mais l’infirmité de ma nature.

Je Vous fais l’offrande anticipée de tous mes combats intérieurs, de toutes mes chutes, de toutes mes larmes, et même de mes éventuelles rébellions involontaires, Vous suppliant de ne pas les regarder comme des refus de ma part, mais comme des misères que je Vous prie de guérir.

Seigneur, mon Dieu, dès aujourd’hui, d’un cœur tranquille et soumis, j’accepte de Votre main le genre de mort qu’il Vous plaira de m’envoyer, avec toutes ses angoisses, toutes ses peines et toutes ses douleurs.

Mon Dieu, faites que je souffre tout ce qu’il Vous plaira, sans jamais me détourner de Vous ; et si je me plains, que ce soit à genoux.

Recevez cette offrande imparfaite, Seigneur, unissez-la aux mérites infinis de Jésus-Christ, et regardez-la par les yeux de Marie, afin qu’elle Vous soit agréable malgré mes défauts.

\conclusion{Fiat Voluntas Tua.}
\conclusion{Ainsi soit-il.}
\culdelampe

\acteverso{Acte d’intention de communion spirituelle perpétuelle}
\usage{quot., le matin ; renouveler par de brèves communions spirituelles dans la journée}
\ouverture{Ô}{}{Jésus,} mon Seigneur et mon Dieu, véritablement présent dans le Très Saint Sacrement de l’autel, je Vous adore avec foi, je Vous aime de tout mon cœur, et je Vous désire plus que tout.

Ne pouvant aujourd’hui Vous recevoir sous les saintes espèces à chaque instant, je forme la résolution sincère de Vous recevoir spirituellement, aussi souvent que ma pensée se tournera vers Vous, avec la même ferveur, le même amour, la même foi qu’à la sainte Table.

Faites, Seigneur, que chaque battement de mon cœur soit comme un acte d’amour eucharistique, une invocation silencieuse, une adoration intérieure, un baiser mystique déposé sur Votre hostie cachée.

Que mon âme soit comme un tabernacle vivant, où Vous puissiez trouver consolation et repos. Venez à moi, ô Jésus, aussi souvent que je soupire vers Vous ; établissez en moi Votre demeure, selon la promesse que Vous ferez à Celui qui garde Votre parole.

Je Vous offre à l’avance toutes les communions spirituelles que je ferai aujourd’hui, dans mes prières, mes silences, mes fatigues et mes élans, et je Vous supplie d’y suppléer ce qui me manque.

Qu’aucune distraction, qu’aucune indifférence, qu’aucun refroidissement ne m’éloigne de Vous ; et si ma nature se montre froide, que mon désir reste brûlant, et que ma fidélité triomphe de ma faiblesse.

Recevez mes actes intérieurs comme autant de communions invisibles ; transformez mes pensées en prières, mes œuvres en offrandes, et mes désirs en feu.

Je désire, Seigneur, que chaque mouvement de mon âme soit une visite au Saint-Sacrement, et que toute ma journée s’écoule comme une longue adoration.

Ô Marie, Vierge fidèle, offrez pour moi à Jésus toutes mes communions spirituelles de ce jour, afin qu’elles soient pures, profondes, persévérantes, et unies à celles que Vous fîtes dans Votre vie terrestre.

\conclusion{Ainsi soit-il.}
\indulg{Acte de communion spirituelle, quelle qu’en soit la formule : 3 ans. — Plénière, aux conditions ordinaires, si l’acte est fait chaque jour pendant un mois. S.~Pénit., 7 mars 1927 et 25 févr. 1933. — Ench.~1952, n\textsuperscript{o}~164.}
\culdelampe

\acteverso{Acte de recommandation des âmes rencontrées}
\usage{quot., le matin}
\ouverture{Ô}{}{Seigneur Jésus-Christ,} Sauveur des âmes, je Vous offre par les mains immaculées de Votre très sainte Mère toutes les âmes que je croiserai aujourd’hui, que je verrai, saluerai, entendrai, ou même seulement apercevrai --- brièvement ou longuement, de loin ou de près --- pour qu’aucune d’entre elles ne m’échappe sans Vous avoir été secrètement consacrée.

Vous les connaissez toutes, Seigneur, dans leur misère et dans leur secret. Je ne demande point de les connaître, ni d’y laisser trace, mais seulement d’être pour elles, selon Votre volonté, un canal muet de grâce, un point d’intercession vivant.

Reprenez-les toutes, ô mon Dieu, par cette offrande matinale : les âmes fidèles, pour qu’elles persévèrent ; les âmes tièdes, pour qu’elles s’éveillent ; les âmes pécheresses, pour qu’elles se convertissent ; les âmes mourantes, pour qu’elles ne désespèrent point. Qu’aucune ne se perde par ma faute, mon oubli ou ma dureté.

Je Vous les confie toutes, ô Marie, Mère du Souverain Prêtre, Refuge des pécheurs, Médiatrice de toutes grâces, Co-Rédemptrice. Parlez-leur à l’instant où je les croise. Agissez en elles par un regard, un silence, une pensée, même inconnue de moi. Disposez de mon attitude, de mon port, de mes paroles, de mon silence, comme d’un voile dont Vous pourriez Vous servir. Prenez-les toutes sous Votre maternelle protection ; et si l’une d’elles devait se damner, qu’au moins je Vous l’aie remise.

Je Vous les offre pour la gloire de Dieu, pour la victoire de l’Église, pour le triomphe de Votre Cœur Immaculé, pour le retour de tous les errants à la Vérité catholique, apostolique et romaine. Que cette offrande soit scellée dans les plaies du Christ, conservée dans Votre Cœur et agréée par le Père éternel.

\conclusion{Ainsi soit-il.}
\culdelampe

%% =====================================================================

\acte{Acte héroïque de charité}
\usage{une fois pour toutes ; à renouveler à volonté, surtout le lundi, jour des défunts, et en novembre}
\begin{notice}
Cet acte de charité parfaite en faveur des âmes du purgatoire, si agréable à Dieu, si utile aux défunts, si profitable pour nous, consiste à donner tous nos mérites satisfactoires à ces saintes âmes. Le pape Benoît XIII a accordé, à ceux qui le font, de très précieuses faveurs. — Bref du 23 août 1728, confirmé par les papes Pie VI et Pie IX. — Voici ces faveurs : 1° indulgence plénière à tout fidèle, à chaque communion et à chaque lundi, où il assiste à la messe pour les défunts ; 2° indulgence plénière aux prêtres, à chacune de leurs messes. Pour gagner ces indulgences, il faut prier quelque peu, aux intentions de notre saint père le pape, dans une visite à l’église. Si l’on entend la messe pour les défunts, le lundi, et si l’on communie, on a deux indulgences plénières à gagner, et deux visites à l’église à faire, en priant aux intentions de notre saint père le pape. On peut prier aux intentions du pape durant la messe elle-même, pour la première visite ; puis, faire une autre visite de quelques minutes, dès la messe achevée. Si l’on ne peut assister à la messe du lundi, on peut gagner cette indulgence plénière à celle du dimanche précédent. 3° En plus, le pouvoir d’appliquer aux défunts toutes les autres indulgences, même celles qui sont accordées seulement pour les vivants ; 4° enfin, tout mérite purement satisfactoire se change en mérites pour le ciel, où nous serons récompensés au centuple, dit saint Ambroise.
\end{notice}
\source{M.-J.-S. Benoît, Livre d’or des âmes du purgatoire, Montréal, 1925 (imprimatur)}
\intertitre{Formule libre de cet acte héroïque}
\ouverture{P}{our}{votre gloire,} ô mon Dieu, pour celle de Marie, la Mère des âmes du Purgatoire, moi, N., je remets entre les mains de cette bonne Mère tous les mérites satisfactoires et tous ceux qui pourront m’être appliqués, en cette vie et en l’autre, afin qu’elle en fasse profiter les âmes qu’elle veut soulager ou délivrer.
\indulg{Plénière, applicable aux seuls défunts : chaque jour où l’on communie, si l’on s’est confessé et si l’on visite une église ou un oratoire public en priant aux intentions du Souverain Pontife ; chaque lundi (ou, si l’on en est empêché, le dimanche suivant), en assistant à la messe pour les défunts, aux conditions ordinaires. Les prêtres jouissent chaque jour de l’autel privilégié personnel. S.~C. des Indulgences, 30 sept. 1852 et 20 nov. 1854. — Ench.~1952, n\textsuperscript{o}~593.}
\culdelampe